\chapter{Selected hints and solutions}
\chaptertag{For checking arguments, not replacing them}

Only selected exercises are treated here. A useful routine is: write a complete attempt, mark the exact step where you are stuck, then consult only the first relevant sentence below.

\section*{Chapter 1}

\textbf{1.} Euclid gives
\[
2024=2(748)+528,\quad748=528+220,\quad528=2(220)+88,
\]
\[
220=2(88)+44,\quad88=2(44),
\]
so the gcd is $44$.

\textbf{2.} The divisions are $527=341+186$, $341=186+155$, $186=155+31$, $155=5\cdot31$. Reverse:
\[
31=186-155=2(186)-341=2(527-341)-341=2\cdot527-3\cdot341.
\]

\textbf{4.} One solution is $35(-5)+22(8)=1$. All solutions are
\[
x=-5+22t,\qquad y=8-35t\qquad(t\in\Z).
\]
To prove completeness, subtract two solutions and use $35(x-x')=-22(y-y')$ with $\gcd(35,22)=1$.

\section*{Chapter 2}

\textbf{3.} Write $c=ar=bs$. Since $a\divides bs$ and $\gcd(a,b)=1$, repeated Euclid's lemma (or B\'ezout) gives $a\divides s$. Thus $s=at$ and $c=abt$.

\textbf{4.} If $\sqrt p=a/b$ in lowest terms, then $a^2=pb^2$. The exponent of $p$ on the left is even; on the right it is odd. Unique factorization forbids this.

\textbf{6.} The primes at most $100$ are
\[
2,3,5,7,11,13,17,19,23,29,31,37,41,43,47,53,59,61,67,71,73,79,83,89,97.
\]
Every composite at most $100$ has a prime factor at most $\sqrt{100}=10$, so
sieving by the primes $2,3,5,7$ crosses out every composite.

\textbf{7.} No: $5!+1=121=11^2$. Euclid's number only needs a prime divisor outside the proposed list; it need not itself be prime.

\textbf{9.} We already know that $a=\sqrt2$ is irrational. By the definition
of the logarithm, $a^b=3$. If $b=p/q$ were rational with $q>0$, then
$(\sqrt2)^{p/q}=3$; raising both sides to $2q$ would give
$2^p=3^{2q}$, contrary to unique factorization. Hence $b$ is irrational.
Unlike the proof by cases, this argument names a particular pair and proves
that it works; it does not leave the witness dependent on an undecided
alternative.

\section*{Chapter 3}

\textbf{1.} Euclid gives
\[
355/113=[3;7,16],
\]
with convergents $3$, $22/7$, and $355/113$.

\textbf{2.} The convergents are
\[
1,\;2,\;\frac32,\;\frac53,\;\frac85,\;\frac{13}{8}.
\]
Numerators and denominators are consecutive Fibonacci numbers; the recurrence with every $a_k=1$ proves it.

\textbf{4.} Since $a_n=(a_n-1)+1/1$ at the bottom level,
\[
[a_0;\dots,a_n]=[a_0;\dots,a_n-1,1].
\]
Requiring the last partial quotient to exceed $1$ selects a unique finite expansion.

\textbf{6.} With $\alpha\approx0.2421988$, the continued fraction begins
\[
 \alpha=[0;4,7,1,3,5,\dots],
\]
so its first four nonzero convergents are
\[
 \frac14,\qquad\frac7{29},\qquad\frac8{33},\qquad\frac{31}{128}.
\]
The Julian calendar uses $1/4$ leap day per year. A century-based cycle has
$q=100Q$, so we instead approximate
\[
 100\alpha\approx24.21988=[24;4,1,1,4,\dots].
\]
Its first nonintegral convergent is $97/4$; dividing the denominator scale by
$100$ gives $97/400$, the Gregorian rate. The corresponding mean years are
\[
 365+\frac14=365.25,
 \qquad 365+\frac{97}{400}=365.2425,
 \qquad 365+\frac{31}{128}=365.2421875.
\]
Relative to the fixed approximation $365.2421988$, their annual errors in days
are about $7.80\times10^{-3}$, $3.01\times10^{-4}$, and
$-1.13\times10^{-5}$, respectively. Thus $31/128$ is numerically better,
whereas $97/400$ fits a convenient rule organized by fours and centuries.

\textbf{8.} Euclid gives $527/341=[1;1,1,5]$, whose convergents are
$1,2,3/2,17/11$. The penultimate convergent is $p_2/q_2=3/2$, so the
remainder formula gives
\[
  \gcd(527,341)=r_2=2\cdot527-3\cdot341=31.
\]

\textbf{9.} The first six convergents of $\sqrt3=[1;\overline{1,2}]$ are
\[
1,\quad 2,\quad \frac53,\quad \frac74,\quad \frac{19}{11},\quad \frac{26}{15}.
\]
They satisfy
\[
1<\frac53<\frac{19}{11}<\sqrt3
<\frac{26}{15}<\frac74<2.
\]
For even $k$, $M_k'(t)<0$, so the finite-tail value lies above
$M_k(\infty)=c_k$; for odd $k$, the derivative is positive and the inequality
reverses. Applying this once to $x$ and once to $c_{k+2}$ gives the two chains.

\textbf{10.} For $k=0$ the formula is just
$D(a_0)=\begin{psmallmatrix}p_0&p_{-1}\\q_0&q_{-1}\end{psmallmatrix}$.
If it holds at $k$, right multiplication by $D(a_{k+1})$ replaces the first
column by $a_{k+1}(p_k,q_k)^T+(p_{k-1},q_{k-1})^T$ and moves the old first
column into second place. These are exactly the recurrences. The determinant
is $(-1)^{k+1}$, so the inverse formula introduces no denominators.

\textbf{11.} For $1/7$, the remainders are
$1,3,2,6,4,5,1,\dots$, producing $0.\overline{142857}$. For $5/12$, they are
$5,2,8,8,\dots$, producing $0.41\overline6$. The first repeated remainder
marks the start and length of the repeating block.

\textbf{15.} Since
\[
D(1)D(2)=\begin{pmatrix}3&1\\2&1\end{pmatrix},
\]
the repeating tail obeys $t=(3t+1)/(2t+1)$. Hence
$2t^2-2t-1=0$, and positivity selects $t=(1+\sqrt3)/2$.

\textbf{16.} For $x=1+\sqrt2$, the conjugate $x'=1-\sqrt2$ lies in
$(-1,0)$; for $x=(1+\sqrt5)/2$, the conjugate is $(1-\sqrt5)/2$, again in
$(-1,0)$. Both numbers exceed $1$, so Galois's criterion applies. Indeed their
expansions are $[\overline2]$ and $[\overline1]$, respectively.

\textbf{17.} Factoring the difference of squares gives
\[
 \abs{\sqrt D-p/q}
 =\frac{\abs{D-p^2/q^2}}{\sqrt D+p/q}
 =\frac{1}{q^2(\sqrt D+p/q)}.
\]
The denominator factor $\sqrt D+p/q$ exceeds $2$, so Legendre's criterion
forces $p/q$ to be a convergent of $\sqrt D$.

\textbf{18.} We have $\sqrt3=[1;\overline{1,2}]$, of even period length $2$;
one period therefore gives the $+1$ solution $2^2-3(1)^2=1$. The first columns
of the first three powers of
$\begin{psmallmatrix}2&3\\1&2\end{psmallmatrix}$ are
\[
 \begin{pmatrix}2\\1\end{pmatrix},\qquad
 \begin{pmatrix}7\\4\end{pmatrix},\qquad
 \begin{pmatrix}26\\15\end{pmatrix},
\]
and each satisfies $p^2-3q^2=1$.

\textbf{19.} A finite computation verifies only the displayed finite averages;
an asymptotic conclusion requires control of every later partial quotient.
Khinchin's theorem excludes only a measure-zero exceptional set and supplies no
known test proving that a named constant such as $\pi$ lies outside it. For $e$,
the repeating pattern of positions gives
\[
 a_{3j-2}=1,\qquad a_{3j-1}=2j,\qquad a_{3j}=1,
\]
so $a_1\cdots a_{3m}=2^m m!$. Taking the $(3m)$th root and applying
$m!\sim\sqrt{2\pi m}(m/e)^m$ yields
\[
 G_{3m}(e)=(2^m m!)^{1/(3m)}\sim(2m/e)^{1/3}\to\infty.
\]
The denominator recurrence gives
$q_{3m}\geq a_1\cdots a_{3m}=2^m m!$, so $q_{3m}^{2/(3m)}\to\infty$ too.
Thus $e$ is provably on the exceptional side of both laws.

\textbf{20.} The two matrices and their maps are
\[
 D(1)D(2)=\begin{pmatrix}3&1\\2&1\end{pmatrix},
 \quad t\longmapsto\frac{3t+1}{2t+1};
 \qquad
 D(2)D(1)=\begin{pmatrix}3&2\\1&1\end{pmatrix},
 \quad t\longmapsto\frac{3t+2}{t+1}.
\]
The rightmost matrix acts first. The fixed-point equations are
$2t^2-2t-1=0$ and $t^2-2t-2=0$, with positive roots
$(1+\sqrt3)/2$ and $1+\sqrt3$, respectively.

\textbf{21.} For $j=1$ and $j=2$, respectively,
\[
 MNM^{-1}
 =\begin{pmatrix}2&3\\1&2\end{pmatrix},
 \qquad
 MNM^{-1}
 =\begin{pmatrix}3&2\\1&1\end{pmatrix}.
\]
The positive fixed points are $\sqrt3$ and $1+\sqrt3$.
Only the second is reduced, since $-\sqrt3<-1$ but
$-1<1-\sqrt3<0$. Its pure period is $[\overline{2,1}]$.

\textbf{22.} A pole $ct+d=0$ would make $t$ rational if $c\ne0$;
if $c=0$, invertibility implies $d\ne0$. If $T_M(t)=r\in\Q$,
then $(a-rc)t=rd-b$. The coefficient $a-rc$ cannot vanish:
together with the equality it would force $ad-bc=0$.
Thus $t=(rd-b)/(a-rc)$ would be rational, a contradiction.
The fixed-point assertion must exclude scalar matrices.
For example, $T_I(t)=t$ for every $t$, including $\sqrt[3]{2}$,
which is not quadratic: $X^3-2$ has no rational root and is an
irreducible cubic over $\Q$.

\textbf{23.} Precisely $1+\sqrt2$, $(1+\sqrt5)/2$, and
$(1+\sqrt3)/2$ are reduced, so those three are purely periodic.
The conjugate of $\sqrt2$ is less than $-1$, and the last two
numbers fail $x>1$. All six are real quadratic irrationals, so
Lagrange's theorem makes all their continued fractions eventually
periodic, including the one with negative initial digit.

\textbf{24.} The successive triples are
\[
 (0,1,3),\ (3,4,1),\ (1,3,1),\ (2,3,1),\
 (1,4,1),\ (3,1,6),\ (3,4,1).
\]
Thus $\sqrt{13}=[3;\overline{1,1,1,1,6}]$.
After the first step, $|P_n|\le3<\sqrt{13}$ and
$0<Q_n\le4<2\sqrt{13}$; each corresponding surd exceeds $1$
and has conjugate in $(-1,0)$.

\textbf{25.} Here $-1/z'=1+\sqrt3$, whose floor is $2$.
The predecessor is $2+1/z=1+\sqrt3$.
The bounds allow $P=-1,0,1$ and $Q=1,2,3$, but the
divisibility condition and both reduced inequalities leave only
$(P,Q)=(1,1),(1,2)$. They represent $1+\sqrt3$ and
$(1+\sqrt3)/2$, and the algorithm interchanges them.
A finite map can have an initial segment before its cycle:
$0\mapsto1$, $1\mapsto1$. Unique reduced predecessors exclude
that behavior in the continued-fraction argument.

\textbf{26.} The period matrix and equation are
\[
 N=\begin{pmatrix}10&3\\7&2\end{pmatrix},
 \qquad 7x^2-8x-3=0,\qquad x=\frac{4+\sqrt{37}}7.
\]
Transposition gives $3z^2-8z-7=0$, with positive root
$z=(4+\sqrt{37})/3$. Finally,
$-1/x'=7/(\sqrt{37}-4)=(4+\sqrt{37})/3$, as predicted.

\textbf{27.} For $D=23$, the triples $(P_n,Q_n,a_n)$ are
\[
 (0,1,4),\ (4,7,1),\ (3,2,3),\ (3,7,1),\ (4,1,8),\ (4,7,1).
\]
Thus $L=4$, the palindrome is $(1,3,1)$, and the final digit is
$8=2a_0$. The first four convergents are $4,5,19/4,24/5$, giving
$24^2-23\cdot5^2=1$. Including the final $8$ instead produces
$211/44$, with remainder $211^2-23\cdot44^2=-7$.

\textbf{28.} The period is $\sqrt6=[2;\overline{2,4}]$.
The four rows $(k,p_k/q_k,Q_{k+1},p_k^2-6q_k^2)$ are
\[
 (0,2/1,2,-2),\ (1,5/2,1,1),\
 (2,22/9,2,-2),\ (3,49/20,1,1).
\]
Their last entries equal $(-1)^{k+1}Q_{k+1}$. Since $L=2$, exactly
the odd-indexed convergents solve $+1$, and none solves $-1$.

\textbf{29.} The least periods for $D=2,3,5,7,13,14$ are, respectively,
\[
 (2),\ (1,2),\ (4),\ (1,1,1,4),\ (1,1,1,1,6),\ (1,2,1,6).
\]
Only $D=2,5,13$ have odd length. Their least positive solutions of
the negative equation are $(1,1),(2,1),(18,5)$.
For $D=3$, reducing $p^2-3q^2=-1$ would give
$p^2\congr2\pmod3$, impossible because a square is $0$ or $1$ modulo $3$.

\textbf{30.} Multiplication by $2-\sqrt3$ sends
\[
 (p,q)\longmapsto(2p-3q,\,2q-p)
\]
and preserves norm $1$. For $q=1$, positivity forces $p=2$, and
the image is $(1,0)$. If $q>1$, then
$\sqrt3<p/q=\sqrt{3+1/q^2}<2$. Hence $2p-3q>0$ and
$0<2q-p<q$. Repeated descent strictly decreases the positive integer
second coordinate until $(2,1)$ and then $(1,0)$ are reached.
Reversing the steps proves the power formula.

\textbf{31.} The state algorithm gives
\[
 \sqrt{61}=[7;\overline{1,4,3,1,2,2,1,3,4,1,14}],\qquad L=11.
\]
The palindrome is the ten-digit block before $14$.
At indices $10$ and $21$, respectively, the convergents give
\[
 29718^2-61\cdot3805^2=-1,
\]
\[
 1766319049^2-61\cdot226153980^2=1.
\]
The second also follows by squaring $29718+3805\sqrt{61}$.
Theorem~\ref{thm:pell-convergents} excludes every smaller positive
solution of either sign.

\textbf{32.} Clear denominators to check the seven equalities
\begin{gather*}
 6W=5B+6Y,\quad20B=9S+20Y,\quad42S=13W+42Y,\\
 12w=7(B+b),\quad20b=9(S+s),\\
 42y=13(W+w),\quad30s=11(Y+y).
\end{gather*}
The displayed table at $t=1$ satisfies each equality and has total
$50\,389\,082$. The first three equations determine all bulls from
$W$. The cow cycle has coefficient
$(7/12)(9/20)(11/30)(13/42)\ne1$, so it determines all cows as well.
Euclid gives gcd $4657$ for the four bull counts and gcd $1$ after
including the cows. A B\'ezout combination of the scaled counts
therefore equals $t$, proving $t\in\Z$ whenever the counts are
integral. At $t=1$, $W+B=4\cdot4\,456\,749$ is not a square,
so the full problem is not solved.

\textbf{33.} If $cr^2=h(h+1)/2$, then
$(2h+1)^2=8cr^2+1$. Conversely, an odd $X>1$ gives the integer
$h=(X-1)/2$ and $h(h+1)/2=(X^2-1)/8=cr^2$.
Here $c=51\,285\,802\,909\,803$.
The identity $8c=4\,729\,494\cdot9314^2$ makes
$(X,r)\mapsto(X,Z)=(X,9314r)$ a bijection onto the positive
solutions of the smaller equation with $9314\divides Z$.
For the reverse map, use $r=Z/9314$; the equation then implies
$X^2\congr1\pmod8$, so $X$ is odd.

\textbf{34.} The recurrence in Section~\ref{sec:archimedes-cattle}
computes the two coefficients of $\varepsilon^j$ modulo $9314$
by induction. Its first vanishing second coefficient occurs at
$j=2329$, with first coefficient $9313\congr-1$.
Because all positive Pell solutions are powers and their second
coordinates increase, this first hit gives the smallest herd.
Write any $j\ge1$ as $j=2329k+r$ with $0\le r<2329$.
Then $\varepsilon^j\congr(-1)^k\varepsilon^r\pmod{9314}$,
so its second coefficient vanishes exactly when $r=0$.
Finally, $T$ is proportional to $Z_j^2$, while
$Z_j=(\varepsilon^j-\varepsilon^{-j})/(2\sqrt{4\,729\,494})$.
Squaring doubles the exponent: $2\cdot2329=4658$.

\section*{Chapter 4}

\textbf{1.} Extended Euclid yields
\[
1=11\cdot101-30\cdot37,
\]
so $37^{-1}\congr-30\congr71\pmod{101}$.

\textbf{2.} Divide $18x\congr30\pmod{42}$ by $6$ while reducing the modulus: $3x\congr5\pmod7$. Since $3^{-1}\congr5$, $x\congr4\pmod7$. Modulo $42$ the six solutions are $4,11,18,25,32,39$.

\textbf{5.} Here $N=504$. One convenient construction is
\[
x=2(72)(4)+3(63)(7)+4(56)(5)=3019\congr499\pmod{504}.
\]
The multipliers $4,7,5$ invert $72,63,56$ modulo $7,8,9$, respectively.

\textbf{8.} The rows $(r,s,t)$ are
\[
 (83,1,0),\ (29,0,1),\ (25,1,-2),\ (4,-1,3),\
 (1,7,-20),\ (0,-29,83).
\]
Thus $1=7\cdot83-20\cdot29$ and $29^{-1}\congr63\pmod{83}$.
Indeed, $29\cdot63=1827=22\cdot83+1$.

\textbf{9.} From $6=42-2\cdot18$, multiplication by $4$ gives
$24=4\cdot42-8\cdot18$. Hence $x\congr-8\congr6\pmod7$;
the residues modulo $42$ are $6,13,20,27,34,41$.
The gcd $6$ does not divide $25$, so the second congruence is impossible.

\textbf{10.} At every component modulus, $e_i$ has value $0$ or $1$,
exactly one selector has value $1$, and distinct selectors cannot both
have value $1$. CRT transfers the three identities back modulo $N$.
Consequently
$(\sum a_i e_i)(\sum b_i e_i)=\sum a_i b_i e_i$ modulo $N$;
addition is componentwise as well.

\textbf{11.} Put $x=2+6t$. Then $6t\congr3\pmod9$, or
$2t\congr1\pmod3$, so $t\congr2\pmod3$ and
$x\congr14\pmod{18}$. The second system prescribes different residues
modulo $\gcd(6,9)=3$, so it is incompatible.

\textbf{12.} The selectors at $-1,0,2$ give
\[
 f(X)=2\frac{X(X-2)}3-\frac{(X+1)(X-2)}2
       +3\frac{X(X+1)}6
     =\frac{2X^2-X+3}{3}.
\]
All solutions are $f(X)+(X+1)X(X-2)h(X)$ with $h\in\Q[X]$.

\textbf{13.} The denominators at $0,1,3$ are $3,5,6$ modulo $7$.
Their inverse certificates are
\[
 3\cdot5-7\cdot2=1,\quad 5\cdot3-7\cdot2=1,\quad
 6\cdot6-7\cdot5=1.
\]
Thus
\[
 f(X)=5(X-1)(X-3)+4\cdot3X(X-3)+2\cdot6X(X-1)
     =X^2+2X+1\quad\text{in }\F_7[X].
\]
Substitution gives the required values $1,4,2$.

\textbf{14.} In $f(2)-f(0)$ every term has a factor $2$, so this
difference cannot be $1$ modulo $8$. The coefficient ring is not a
field; specifically, the difference between the two interpolation
points is a nonunit.

\textbf{15.} Modulo $2$, the orbit is
$(0,1),(1,1),(1,0),(0,1)$, giving period $3$.
Modulo $3$, it is
\[
 (0,1),(1,1),(1,2),(2,0),(0,2),(2,2),(2,1),(1,0),(0,1),
\]
giving period $8$. At the first positive zero, $n=4$, the next residue
is $2$, not $1$, so the recurrence has not restarted.

\textbf{16.} The periods are $\operatorname{lcm}(3,8)=24$ and
$\operatorname{lcm}(8,20)=40$. A simultaneous return occurs exactly at
an index divisible by both component periods. For modulus $15$,
the product $8\cdot20=160$ is a period but is four times the least one.

\textbf{17.} The orbit is $(1,0)\mapsto(0,0)\mapsto(0,0)$.
Modulo $2$ the rule is $(a,b)\mapsto(b,b)$, which forgets $a$;
a state therefore does not determine a unique predecessor.

\textbf{18.} The squares modulo $7$ are $0,1,2,4$. A root $x$ would
satisfy $(2x-1)^2=5$, which is impossible. Binet's formula requires
the missing roots (equivalently, a square root of $5$); the extension
in Appendix~\ref{app:finite-fields} supplies arithmetic for them.

\section*{Chapter 5}

\textbf{1.} The binary length-$3$ words form four orbits:
\[
\{000\},\quad\{111\},\quad\{001,010,100\},\quad\{011,110,101\}.
\]

\textbf{3.} If $ik\congr jk\pmod p$, then $p\divides k(i-j)$. As $p\ndivides k$, Euclid's lemma gives $p\divides i-j$. For $0\le i,j<p$, this forces $i=j$.

\textbf{5.} Orbit sizes divide $6$, and binary examples realize $1,2,3,6$: $000000$, $010101$, $001001$, and $000001$.

\textbf{6.} Fermat gives $(a+b)^p=a+b=a^p+b^p$ in $\F_p$.
Applying it repeatedly proves the corresponding identity for any finite
sum. Fermat's statement about residues does not assert $u^p=u$ for
new elements of an extension; the polynomial identity from
Chapter~\ref{chap:binomial} proves additivity there instead.

\section*{Chapter 6}

\textbf{1.} $(\Z/12\Z)^\times=\{1,5,7,11\}$. Every nonidentity element squares to $1$, so their orders are $2$; the identity has order $1$.

\textbf{4.} If $aH\cap bH$ contains $x$, write $x=ah_1=bh_2$. Then $a=b h_2h_1^{-1}$, so every $ah\in aH$ equals $b(h_2h_1^{-1}h)\in bH$. Thus $aH\subseteq bH$; symmetry gives equality.

\textbf{5.} For $g\ne e$, the order of $g$ divides the prime $\abs G$, so it is $\abs G$. Hence $G=\angles g$.

\section*{Chapter 7}

\textbf{2.} Legendre gives
\[
v_2(100!)=50+25+12+6+3+1=97,
\]
\[
v_5(100!)=20+4=24.
\]
Each trailing zero consumes one $2$ and one $5$, so there are $24$.

\textbf{4.} Choose $n$ objects from two disjoint groups of $n$. If $k$ come from the first group, then $n-k$ come from the second. Summing over $k$ gives
\[
\sum_k\binom nk\binom n{n-k}=\binom{2n}{n}.
\]

\textbf{7.} Legendre gives
\[
v_p\binom{2n}{n}=\sum_{j\ge1}\left(\floor{\frac{2n}{p^j}}-2\floor{\frac n{p^j}}\right).
\]
Each summand is $0$ or $1$, and there are at most $\floor{\log_p(2n)}$ nonzero candidates. This is the exponent of $p$ in the lcm, so the valuation of the binomial coefficient is no larger for every prime $p$.

\textbf{8.} Fermat makes $X^p-X$ zero at every point of $\F_p$,
but its leading coefficient is $1$, so it is not the zero polynomial.
By contrast, expansion shows that every coefficient of
$(X+Y)^p-X^p-Y^p$ vanishes modulo $p$. This is a polynomial identity
and remains valid upon substitution in any field containing $\F_p$.

\section*{Chapter 8}

\textbf{2.} The sum of the $2n+1$ coefficients in row $2n$ is $4^n$. The central coefficient is largest, so it is at least the average $4^n/(2n+1)$ and strictly below the whole sum.

\textbf{3.} For $n<p\le2n$, the numerator $(n+1)\cdots(2n)$ contains exactly one multiple of $p$, while $n!$ contains none. Because $p^2>2n$ for $p>n\ge1$, the valuation is exactly one.

\textbf{7 (hint).} Insert $x=p_n$ into both Chebyshev bounds to compare $n$ with $p_n/\log p_n$. First obtain coarse polynomial bounds on $p_n$ so that $\log p_n$ can be compared with $\log n$, then substitute back.

\section*{Appendix B: continued-fraction arithmetic}

\textbf{1.} With $t_2\in[1,\infty]$,
\[
 x=2+\frac1{3+1/t_2}.
\]
This expression increases with $t_2$.  At $t_2=1$ it is $9/4$; as
$t_2\to\infty$ it approaches $7/3$.  For an infinite continued fraction the
second endpoint is not attained, so $x\in[9/4,7/3)$.

\textbf{2.} The denominator is positive for $t\geq1$, and
\[
 F'(t)=\frac{22\cdot2-3\cdot14}{(14t+2)^2}
       =\frac2{(14t+2)^2}>0.
\]
Thus the range is determined by the endpoint values
$F(1)=25/16$ and $F(\infty)=11/7$.  The whole interval lies in $[1,2)$, so
the certified output digit is $1$.

\textbf{3.} Substituting $t=m+1/u$ into
$F_M(t)=(At+B)/(Ct+D)$ and clearing the factor $u$ gives
\[
 \frac{(Am+B)u+A}{(Cm+D)u+C},
\]
which is represented by $MD(m)$.  After emitting $k$, the unread output is
$1/(F_M-k)$, represented by
$\begin{psmallmatrix}0&1\\1&-k\end{psmallmatrix}M$.

\textbf{4.} From $M=\begin{psmallmatrix}2&0\\0&1\end{psmallmatrix}$,
absorption gives
\[
 MD(15)=\begin{pmatrix}30&2\\1&0\end{pmatrix},
 \qquad
 MD(15)D(1)=\begin{pmatrix}32&30\\1&1\end{pmatrix}.
\]
The latter map has range $[31,32)$ on $[1,\infty)$, so the next output digit
is $31$.

\textbf{5.} For each fixed $v$, the map in $u$ is linear fractional and hence
monotone unless it is constant; its extrema on a horizontal slice occur on the
two vertical edges.  On either edge, the remaining function of $v$ is again
linear fractional, so its extrema occur at the two endpoints.  Therefore all
global extrema are corner values.  If the denominator vanishes inside the
rectangle, the map is not continuous there and can become unbounded, so corner
values alone give no range certificate.

\textbf{7.} Treating the two copies as independent discards the relation
$x=y$.  Every finite continued-fraction prefix gives a nontrivial interval
containing $\sqrt2$, so the product box permits values slightly below and
slightly above $2$.  Since the floor jumps at $2$, it is never certified from
such boxes.  Retaining the dependency $x=y$ changes the expression to $x^2$;
using the known equation $x^2=2$ resolves it immediately.  Without that extra
algebraic relation, an interval-valued output can still certify arbitrarily
small numerical error.

\section*{Appendix C: Venus and astronomical cycles}

\textbf{1.} $E$ is Earth's Sun-centered lap time and $V$ is Venus's
Sun-centered lap time.  The quantity $S$ measures something different: the
waiting time between successive Earth--Venus alignments.

\textbf{2.} Venus and Earth move through $1/V$ and $1/E$ laps per day.
Their relative phase therefore increases at $1/V-1/E$ revolutions per day, so
one relative revolution takes $S$ days with $1/S=1/V-1/E$.  If $r=E/V$,
then
\[
  \frac ES=\frac EV-1=r-1,
  \qquad \frac SE=\frac1{r-1}.
\]
Substitution gives $T(13/8)=1/(13/8-1)=8/5$.

\textbf{3.} Repeated addition of $3/5$ modulo $1$ gives
\[
  0,\ \frac35,\ \frac65\congr\frac15,\
  \frac45,\ \frac75\congr\frac25,
\]
then returns to $0$.  Each step skips one neighboring vertex of the regular
pentagon, so the chronological edges are the diagonals and form a pentagram.

\textbf{5.} The factorizations are
\[
  2920=2^3\cdot5\cdot73,
  \qquad 260=2^2\cdot5\cdot13.
\]
Thus the gcd is $20$ and the lcm is
$2^3\cdot5\cdot13\cdot73=37960$.  Division by $584,365,260,$ and $2920$
gives the multipliers $65,104,146,$ and $13$.

\textbf{6.} The Fibonacci recurrence gives
\[
 T\!\left(\frac{F_{n+1}}{F_n}\right)
 =\frac{1}{F_{n+1}/F_n-1}
 =\frac{F_n}{F_{n+1}-F_n}
 =\frac{F_n}{F_{n-1}}.
\]
The positive solution of $r=1/(r-1)$ is
$r=(1+\sqrt5)/2=\varphi$.  Its identity
$\varphi=1+1/\varphi$ is exactly the infinite continued fraction
$[1;\overline1]$.

\textbf{7.} The nominal count gains $0.08$ day per cycle, hence $0.4$ day
after $5$ cycles, $5.2$ days after $65$, and $52$ days after $650$.

\section*{Appendix \ref{app:finite-fields}: extension arithmetic and periods}

\textbf{1.} The product is $13+7s=6$. Since
$1-5\cdot2^2=-19=2$ in $\F_7$, its inverse is
$4(1-2s)=4+6s$. Multiplication gives
$(1+2s)(4+6s)=64+14s=1$.

\textbf{2.} The elements are $0,1,t,t+1$. The polynomial
$X^2+X+1$ has no root in $\F_2$ and is irreducible. Directly,
$t(t+1)=t^2+t=1$, so the nonzero inverses are $1,t+1,t$,
respectively. Squaring fixes $0,1$ and exchanges $t,t+1$.

\textbf{3.} Write $u=a+bs$. Then $u^7=a-bs$, which equals $u$
exactly when $2b=0$, hence $b=0$. Applying the map twice gives
$u^{49}=u$. The example $s^7=-s\ne s$ exhibits its nonidentity action.

\textbf{4.} Both $6$ and $24$ satisfy $X^2-X-1=0$ modulo $29$.
We have $6^7=-1$, $6^{14}=1$, and $6^2=7\ne1$, so $6$ has
order $14$. Also $24^7=1$ and $24\ne1$, so $24$ has order $7$.
Their least common multiple is $14$, agreeing with factor removal.

\textbf{5.} Frobenius interchanges the roots, so
$\alpha^{p+1}=\beta^{p+1}=-1$. Binet gives $F_{p+1}=0$ and
$F_{p+2}=(-\alpha+\beta)/(\alpha-\beta)=-1$.
For $p=3$ this is precisely the state $(0,2)$ at index $4$;
the full return occurs at index $8$.

\textbf{6 (check values).} For primes $3,7,11,13,29$, valid starting
multiples are $8,16,10,28,28$ and the respective least periods are
$8,16,10,28,14$. The exceptional primes $2,5$ have periods $3,20$.
The modular arithmetic notebook gives implementations of direct pair
iteration, fast doubling, and factor removal for comparison.
